Write \(R_{\rm full}\) and \(R_{\rm hit}\) for the two displayed minimax risks.  The target-hit record is a measurable image of the full record. hence every rule for the hit experiment lifts to a full-data rule.  Consequently
\[
 R_{\rm full}\leq R_{\rm hit}. \tag{1}
\]
For the converse, parameterize the observed conditional law by the assignment-specific marginals \(F_z=\mathcal L_P((Y_j(z))_{j=1}^n)\).  Any collection \((F_z:z\in\mathcal Z_n)\) is induced by the legal schedule law \(\bigotimes_{z\in\mathcal Z_n}F_z\).  In particular, fix once and for all known nontarget marginals \(F_z^0\), and let \(\mathcal P_0\) consist of product schedule laws for which \(F_z=F_z^0\) off \(\bigcup_k\mathcal S_k\), while the target marginals are arbitrary.  This restriction does not restrict \(U\), the target-hit law, or the set of possible collections of target marginals.

Under \(\mathcal P_0\), conditional on the reduced record, every omitted nontarget outcome vector has its known law \(F_{Z_c}^0\), independently across clusters.  Therefore, from any full-data randomized rule \(\delta\), construct a reduced-data randomized rule \(\delta^0\) by drawing these omitted vectors from their known conditional laws and applying \(\delta\).  For every \(P\in\mathcal P_0\), the joint law of the completed record and \(U(P)\) is exactly its full-experiment law, so
\[
 \mathbb E_P\mathcal L(U,\delta^0)=\mathbb E_P\mathcal L(U,\delta).
\]
It follows, using first the restriction \(\mathcal P_0\subseteq\mathcal P\), then the fact that its target marginals exhaust the reduced parameter space, that
\[
 \sup_{P}\mathbb E_P\mathcal L(U,\delta)
 \geq \sup_{P\in\mathcal P_0}\mathbb E_P\mathcal L(U,\delta)
 =\sup_P\mathbb E_P\mathcal L(U,\delta^0)
 \geq R_{\rm hit}.
\]
Taking the infimum over \(\delta\) and combining with (1) proves equality.

For target records, the preceding factorization gives the conditional kernel
\[
 \Pr(L_c=\ell\mid Z_c=z)=p_\ell B_{k\ell}/q_k,\qquad z\in\mathcal S_k,
\]
which is known and independent of every \(F_z\).  A rule observing the label can therefore be reproduced after label deletion by simulating from this kernel.  Forgetting the label gives the reverse comparison, so the same two inequalities prove the final assertion.